\section{Background}
\label{sec:background}

An inertial navigation system integrates angular rate and specific force to propagate position, velocity and attitude, and its error grows without bound unless an external measurement corrects it \citep{groves}. In most vehicles that correction is GNSS. Interference, spoofing and plain unavailability can remove or corrupt it \citep{9112619}, which has driven interest in alternative position sources such as gravity and magnetic anomaly maps \citep{canciani2016absolute,canciani2017airborne,gravnav}. Developing and comparing such methods requires recordings in which the inertial error has time to grow. With consumer MEMS sensors that means drives of tens of minutes to hours, and for map-based aiding it also means distances long enough to cross map features. The public grids used here have cells of 1 arc-minute for the free-air gravity anomaly and 3 arc-minutes for the magnetic anomaly, about 1.9 and 5.6~km \citep{SANDWELL20211059,wdmam}.

Existing datasets fit this need poorly. Vehicle and robot datasets such as KITTI \citep{kitti}, Oxford RobotCar \citep{oxford-robotcar}, the Ford campus set \citep{ford-campus}, the DARPA Urban Challenge set \citep{mit-darpa}, KAIST Urban \citep{kaist-urban} and NCLT \citep{nclt} pair navigation-grade inertial or GNSS/INS references with cameras and lidar, and they are recorded on campuses or in cities. Aerial and handheld visual--inertial sets \citep{burri2016euroc,tum-vi,uzh-fpv,nguyen2022ntu,newer-college} last minutes. Smartphone datasets exist, but they target other problems: raw GNSS measurements for precise positioning \citep{fu2020android}, or pedestrian inertial odometry \citep{herath2020ronin}. Inertial datasets for autonomous platforms \citep{ansfl,yampolsky2024multiple} are built around dedicated inertial units. To our knowledge, no public dataset combines hours-long road drives, consumer phone inertial sensors from several manufacturers, barometer and magnetometer streams, and a ready-made way to impose controlled GNSS degradation.

An earlier version of MEMS-Nav was released as a preprint and a GitHub repository with 17 trajectories at about 1~Hz \citep{memsnav-preprint}, and parts of it were used to evaluate anomaly-aided filters \citep{anom-ukf,anom-pf}. This article supersedes that release. It adds the raw recordings at their delivered rates, \numTrajectories{} trajectories in place of 17, a 10~Hz processed form, per-device and per-channel documentation, map tiles for every trajectory, reference baselines under a specified degradation scenario, and a characterization of the phone channels that explains why they bring no measurable benefit when used as anomaly measurements \citep{brodovsky2026dissertation}.
